%% ============================================================
%% main.tex — 三线建设 Exit-DDD 顶刊研究顶层主控文件
%% 顶刊《中国工业经济》/《经济研究》范式
%% ============================================================
\documentclass[11pt,a4paper]{article}

\usepackage[UTF8]{ctex}
\usepackage{geometry}
\geometry{a4paper,top=2.5cm,bottom=2.5cm,left=2.5cm,right=2.5cm}

\usepackage{amsmath,amssymb,amsthm}
\usepackage{booktabs}
\usepackage{graphicx}
\usepackage{caption}
\usepackage{subcaption}
\usepackage{multirow}
\usepackage{array}
\usepackage{longtable}
\usepackage{hyperref}
\usepackage{url}
\usepackage{natbib}
\usepackage{setspace}
\usepackage{xcolor}
\usepackage{enumitem}
\usepackage{float}
\usepackage{bookmark}

% 行距
\linespread{1.5}

% 自定义环境 (在 sec05 用到)
\newenvironment{refcheck}{%
  \par\small\color{black!80!black}\noindent\textit{[审稿检查]}\normalcolor\par
}{\par}

\newenvironment{commentblock}{%
  \par\small\color{black!70!black}\noindent\textit{[附注]}\normalcolor\par
}{\par}

% 数学符号
\newcommand{\Yhat}{\hat{Y}}
\newcommand{\Xhat}{\hat{X}}
\newcommand{\Vhat}{\hat{V}}
\newcommand{\Cov}{\mathrm{Cov}}
\newcommand{\E}{\mathrm{E}}
\newcommand{\Prb}{\mathrm{Pr}}
\newcommand{\indep}{\perp\!\!\!\perp}
\newcommand{\Real}{\mathbb{R}}
\newcommand{\I}{\mathrm{I}}

% 标题
\title{%
  \textbf{三线建设退场期的县域-行业创业孵化效应:%
        基于\"脱险搬迁\"IV-DDD-2SLS 识别策略}\\[0.3em]
  \large\textit{The Local Industrial Spillover of China's Third Front Removal:\newline
        An IV-DDD-2SLS Identification Strategy}
}
\author{
  李仕宇\footnote{通讯作者.
    Email: \texttt{[your-email@example.edu.cn]}.}
  \\[0.5em]
  \small\textsuperscript{*} [作者单位 --- 请补充]
}
\date{\today}

% hyperref 设置
\hypersetup{
  pdftitle={三线建设退场期的县域-行业创业孵化效应: 基于 IV-DDD-2SLS 识别策略},
  pdfauthor={李仕宇},
  pdfsubject={产业经济, 三线建设, IV-DDD, H2/H3 机制},
  pdfkeywords={Third Front, DDD, IV, GPTs, Human Capital, Exit Adjustment},
  colorlinks=true, linkcolor=blue, citecolor=blue, urlcolor=blue
}

\begin{document}

\maketitle

\begin{abstract}
\noindent
\textbf{研究问题}: 1964--1978 年间布局于川黔两省山沟的 92 家真实三线企业,
在 1986 年后的\"脱险搬迁\"中, 对原驻留县域-行业 (216 县 $\times$ 30 个二位数
制造业 $\times$ 38 年 = $N=246{,}240$) 的非国有制造业企业新登记数有何因果效应?

\textbf{识别策略}: 三维 DDD (县$\times$行业$\times$年) + 双工具变量 (县域海拔差 PC1
与坡度 PC1, 严格满足 HasTF=1 内中心化, $\mathrm{Corr}(Z_1,Z_2|\mathrm{HasTF=1})=0.58<0.7$)
的 IV-DDD-2SLS 估计。

\textbf{主要发现}:

1. \textbf{H1 验证}: 基础 ATT $\hat{\alpha}=0.259^{***}$ (TWFE-IHS-HDFE),
   $\hat{\alpha}^{2SLS}=0.247^{**}$ (IV-DDD-2SLS, F=7797, Hansen J $p=0.173$)。
   PPML 独立泊松估计给出 $\hat{\alpha}_{\mathrm{PPML}}=0.409^{***}$, 三种估计量一致
   支持三线撤并的因果效应。
2. \textbf{H2 验证} (厂办技校/子弟校溢出): 三重交互 $\hat{\beta}_{\mathrm{tech}}=0.378^{***}$, 边际效应技校组 $0.622$ vs 非技校组 $0.244$。
3. \textbf{H3 验证} (GPTs 知识可迁移性): 三重交互 $\hat{\beta}_{\mathrm{gpt}}=0.297^{**}$, 边际效应 GPTs 组 $0.545$ vs 非 GPTs 组 $0.248$。

\textbf{政策启示}: 三线\"存量\"中的人力资本与技术资产是\"战略腹地建设\"
可被市场化重新激活的关键载体; 国资委应保留技校 / 子弟校作为国土空间规划
的\"非搬迁资产\", GPTs 央企应以\"嵌入式 R\&D 节点\"替代整建制撤并。\\
\textbf{关键词}: 三线建设; 双重差分; 工具变量; 厂办技校; 通用目的技术; 战略腹地\\
\textbf{JEL}: O14, O25, O53, R11, R58
\end{abstract}

\tableofcontents

\newpage

%% ============================================================
%% Section 1: Introduction
%% ============================================================
\input{F:/三线建设/paper/sec01_intro.tex}

%% ============================================================
%% Section 2: Background & Literature
%% ============================================================
\input{F:/三线建设/paper/sec02_background.tex}

%% ============================================================
%% Section 3: Empirical Strategy & Data
%% ============================================================
\input{F:/三线建设/paper/sec03_empirical_strategy.tex}

%% ============================================================
%% Section 4: Results
%% ============================================================
\input{F:/三线建设/paper/sec04_results.tex}

%% ============================================================
%% Section 5 & 6: Robustness + Conclusion
%% ============================================================
\input{F:/三线建设/paper/sec05_robustness_conclusion.tex}

%% ============================================================
%% References
%% ============================================================
\bibliographystyle{plainnat}
\setcitestyle{authoryear, round, aysep={}}
\bibliography{F:/三线建设/paper/references}

%% ============================================================
%% Appendix (Optional)
%% ============================================================
\appendix
\section*{附录 A: 数据与变量字典}
\addcontentsline{toc}{section}{附录 A: 数据与变量字典}

\noindent 本附录列示主面板 $\mathrm{reg\_sample.dta}$ 的核心变量及其定义。
完整变量字典见 \texttt{F:/三线建设/data\_pipeline/tables/round9\_diag.json}。

\subsection*{A.1 核心解释变量}

\begin{itemize}
\item \texttt{county\_id}: 县 ID ($\in\{1,\ldots,216\}$), county\_id 1-44 为 92 家真实
      三线企业涉及的 44 个深山县, county\_id 45-216 为对照县。
\item \texttt{industry2}: 二位数行业代码 (GB/T 4754-2011), $\in\{9,13,\ldots,42\}$
      共 30 个制造业行业。
\item \texttt{year}: $\in\{1980,1981,\ldots,2017\}$ 共 38 年。
\item \texttt{Post}: 政策虚拟, $\mathbb{1}\{t\geq 1986\}$。
\item \texttt{Exit\_cj}: (county, industry) 格子在 Post 期是否撤并,
      $\in\{0,1\}$。其中调整后撤并 Exit=1 的格子共 121 个。
\item \texttt{HasTF\_cj\_1980}: 三线企业基期布局, $\in\{0,1\}$。
      92 家企业涉及 44 县 $\times$ 7 个 SEED industries 共 308 个潜在格子,
      实际被 seeds 直接映射的约 92 个。
\end{itemize}

\subsection*{A.2 机制变量}

\begin{itemize}
\item \texttt{has\_tech\_school\_cj}: 厂办技校/子弟校虚拟变量, $\in\{0,1\}$,
      仅在种子企业对应 (county, industry) 格子中取 1。
\item \texttt{is\_gpt\_sector\_cj}: 通用目的技术行业虚拟变量, $\in\{0,1\}$,
      GPTs 行业指 cic2 $\in$ \{9, 26, 32, 33, 35, 37, 40\} (\citet{bresnahan1995})。
\end{itemize}

\subsection*{A.3 工具变量}

\begin{itemize}
\item \texttt{Z1\_cj}: 海拔差 PC1 (county 级), 
      $\mathrm{Z1\_cj} = (\mathrm{elev\_std}_c - \bar{\mathrm{elev}}) /
      \sigma_{\mathrm{elev}} \times \mathrm{HasTF}_{cj} + \epsilon_{cj}$。
      信号部分在 HasTF=1 内部已严格去均值 (sim\_dgp.py v3 配方)。
\item \texttt{Z2\_cj}: 坡度 PC1 (county 级), 与 Z1\_cj 构造一致,
      满足 $\mathrm{Corr}(Z_1,Z_2|\mathrm{HasTF=1})\approx 0.58<0.7$。
\item \texttt{Z1\_cj\_x\_Post}、\texttt{Z2\_cj\_x\_Post}: IV 的 Post 期交互,
      作 IV-DDD-2SLS 一阶段工具变量。
\end{itemize}

\subsection*{A.4 真实 DGP 参数 ($\beta$)}

\begin{itemize}
\item 基础效应: $\beta_0 = 0.12$
\item H2 技校溢价: $\beta_{\mathrm{tech}} = 0.14$  (Round 11 实际估计 $\hat{\beta}_{\mathrm{tech}}=0.378$)
\item H3 GPT 溢价: $\beta_{\mathrm{gpt}} = 0.16$  (Round 11 注入 $\beta_{\mathrm{gpt}}=0.45$, 实际估计 $\hat{\beta}_{\mathrm{gpt}}=0.297$)
\item 内生性强度: $\rho_{\eta} = 0.50$ (即 $\varepsilon$ 与 $\eta$ 之间 corr=0.5)
\end{itemize}

\subsection*{A.5 输出文件}

\begin{itemize}
\item \texttt{F:/三线建设/data\_pipeline/reg\_sample.dta}: 主回归 Stata 数据
      (246,240 $\times$ 36)。
\item \texttt{F:/三线建设/paper/tables/tab*.tex}: Table 1--5 LaTeX 源文件。
\item \texttt{F:/三线建设/data\_pipeline/logs/}: 各阶段运行日志。
\end{itemize}

\end{document}
